% SGA 3，第 VII_A 讲， 英文重建稿——第 1 轮。
% 范围：权威底本局部页 33--37, 段落 5.3--5.3.4,
% 止于其前： 5.4 局部页 37.
% 合并阅读本页码： 486--490. 再版印刷页码： 475--479.
% 本范围内原版页边定位： 448--452.
% 权威底本： Polo--Gille Exp7A-13oct24.pdf,
% SHA-256 D46DB8C79D0385643B6EED4B96900613248375F9B73B4202478A4F57FD3E8522.
% 曾参考 Jacob Reinhold 的 CC BY 4.0 Markdown 作为比对谱系。

\medskip
\noindent\textbf{5.3.}
\phantomsection
\label{sga3:VIIA:5.3}

% 权威底本： Exp7A 局部页 33; 合并本页 486; 印刷页 475.
现在我们研究函子
\[
  A\longmapsto A_{\mathrm{Lie}}
\]
（见 \SGARefLink{sga3:VIIB:5.2:theorem}{5.2}）的左伴随。设 \(\mathfrak g\) 是特征为 \(p\) 的环 \(R\) 上的受限 \(p\)-Lie 代数，设
\(U(\mathfrak g)\) 为 \(\mathfrak g\) 的底层 Lie 代数的包络代数（参见 \textup{[BLie]}, I, \S2.1），并设
\[
  i_{\mathfrak g}:\mathfrak g\longrightarrow U(\mathfrak g)
\]
（或简记为 \(i\)）为典范映射。

设 \(A\) 是含幺结合 \(R\)-代数。对于每个 Lie 代数同态
\[
  \varphi:\mathfrak g\longrightarrow A_{\mathrm{Lie}},
\]
存在唯一的含幺 \(R\)-代数同态
\[
  \psi:U(\mathfrak g)\longrightarrow A
\]
使得 \(\psi\circ i=\varphi\)。此外，\(\varphi\) 是受限 \(p\)-Lie 代数同态，当且仅当 \(\psi\) 在所有下列元素上为零：
\[
  i(x)^p-i\bigl(x^{(p)}\bigr),
  \qquad x\in\mathfrak g.
\]

\noindent\textbf{定义。}
将 \(U(\mathfrak g)\) 对由元素
\[
  i(x)^p-i\bigl(x^{(p)}\bigr)
\]
生成的双边理想所取的商记为 \(U_p^R(\mathfrak g)\)，或简记为
\(U_p(\mathfrak g)\)。将复合映射记为 \(j_{\mathfrak g}\)，或简记为 \(j\)：
\[
  \mathfrak g
  \xrightarrow{i}
  U(\mathfrak g)
  \longrightarrow
  U_p(\mathfrak g).
\]
代数 \(U_p(\mathfrak g)\) 称为 \(\mathfrak g\) 的\emph{受限包络代数}。

上述讨论给出以下结果。

\medskip
\noindent\textbf{命题。}
\textit{对于每个含幺结合 \(R\)-代数 \(A\) 以及每个受限 \(p\)-Lie 代数同态
\[
  \varphi:\mathfrak g\longrightarrow A_{\mathrm{Lie}},
\]
存在唯一的含幺代数同态
\[
  \psi:U_p(\mathfrak g)\longrightarrow A
\]
使得 \(\psi\circ j=\varphi\)。换言之，函子
\(\mathfrak g\mapsto U_p(\mathfrak g)\) 是遗忘函子 \(A\mapsto A_{\mathrm{Lie}}\) 的左伴随。}

\medskip
\noindent\textbf{5.3.1.}
\phantomsection
\label{sga3:VIIA:5.3.1}

% 权威底本： Exp7A 局部页 34; 合并本页 487; 印刷页 476.
沿用 \SGARefLink{sga3:VIIA:5.3}{5.3} 的记号，令
\[
  \beta(x)=i(x)^p-i\bigl(x^{(p)}\bigr).
\]
对于每个 \(y\in\mathfrak g\)， \SGARefLink{sga3:VIIA:5.1}{公式 5.1}~\SGARefLink{sga3:vii:paragraph:5-1:equation:i}{(i)} 和 \SGARefLink{sga3:VIIB:5.2:theorem}{5.2}~\SGARefLink{sga3:vii:paragraph:5-1:equation:i}{(i)}
给出
\[
\begin{aligned}
  \beta(x)i(y)
  &=
  i(y)\beta(x)+[\beta(x),i(y)]\\
  &=
  i(y)\beta(x)
  +(\operatorname{ad}i(x))^p i(y)
  -i\bigl((\operatorname{ad}x)^p y\bigr)\\
  &=
  i(y)\beta(x).
\end{aligned}
\]
因此 \(\beta(x)\) 属于 \(U(\mathfrak g)\) 的中心；特别地，由元素
\(\beta(x)\) 生成的左理想已经是双边理想。

另一方面，
\[
  \beta(\lambda x)=\lambda^p\beta(x)
  \qquad(\lambda\in R),
\]
且 \SGARefLink{sga3:VIIA:5.1}{公式 5.1}~\SGARefLink{sga3:vii:paragraph:5-1:equation:ii}{(ii)} 和 \SGARefLink{sga3:VIIA:5.2}{5.2}~\SGARefLink{sga3:vii:paragraph:5-1:equation:ii}{(ii)} 给出
\[
  \beta(x+y)=\beta(x)+\beta(y)
  \qquad(x,y\in\mathfrak g).
\]
因而，若 \((x_\alpha)\) 是 \(R\)-模 \(\mathfrak g\) 的一族生成元，则由所有
\(\beta(x)\) 生成的左理想已经由元素 \(\beta(x_\alpha)\) 生成。

\medskip
\noindent\textbf{5.3.2. 命题。}\footnote{本段中，编者调整了顺序，先陈述结果，再给出详细证明。}
\phantomsection
\label{sga3:VIIA:5.3.2}
\textit{设 \(\mathfrak g\) 是一个 \(R\)-Lie 代数，其底层 \(R\)-模为自由模，基为 \((x_\alpha)\)。则
\(\mathfrak g\) 上的受限 \(p\)-Lie 代数结构与满足下式的 \(\mathfrak g\) 中元素族 \((y_\alpha)\) 一一对应：
\[
  \operatorname{ad}y_\alpha
  =
  (\operatorname{ad}x_\alpha)^p.
\]}

事实上，设 \(\mathfrak g\) 配备受限 \(p\)-Lie 代数结构 \(x\mapsto x^{(p)}\)。由 \SGARefLink{sga3:VIIB:5.2:theorem}{5.2}~\SGARefLink{sga3:vii:paragraph:5-1:equation:i}{(i)},
\SGARefLink{sga3:VIIB:0}{(0)}, 及 \SGARefLink{sga3:vii:paragraph:5-1:equation:ii}{(ii)}，元素
\[
  y_\alpha=x_\alpha^{(p)}
\]
满足
\[
  \operatorname{ad}y_\alpha
  =
  (\operatorname{ad}x_\alpha)^p,
\]
并且它们确定了该受限 \(p\)-Lie 代数结构。

下面证明逆命题。由于 \(\mathfrak g\) 是自由 \(R\)-模，由 Poincar\'e--Birkhoff--Witt 定理（参见 \textup{[BLie]}, I, \S2.7），典范映射
\[
  i:\mathfrak g\longrightarrow U(\mathfrak g)
\]
是单射。因此可将 \(\mathfrak g\) 等同于 \(U(\mathfrak g)\) 的一个 \(R\)-子模。设 \((y_\alpha)\) 是 \(\mathfrak g\) 中满足下式的一族元素：
\[
  \operatorname{ad}y_\alpha
  =
  (\operatorname{ad}x_\alpha)^p.
\]
令
\[
  \pi:R\longrightarrow R,
  \qquad
  r\longmapsto r^p,
\]
并令 \(\mathfrak g\otimes_\pi R\) 为沿 \(\pi\) 作标量扩张所得的 \(R\)-Lie 代数。\footnote{也就是说，
\(xr\otimes_\pi1=x\otimes_\pi r^p\) 且
\(r\cdot(x\otimes_\pi1)=x\otimes_\pi r\)，其中 \(x\in\mathfrak g\)，\(r\in R\)。}

存在一个 \(R\)-线性映射
\[
  \gamma:\mathfrak g\otimes_\pi R
  \longrightarrow
  U(\mathfrak g)
\]
将 \(x_\alpha\otimes_\pi1\) 映至
\(x_\alpha^p-y_\alpha\)。此外，对于每个 \(x\in\mathfrak g\)，
\[
  (\operatorname{ad}x_\alpha^p)(x)
  =
  (\operatorname{ad}x_\alpha)^p(x)
  =
  (\operatorname{ad}y_\alpha)(x).
\]
由此可知，\(\gamma\) 将 \(\mathfrak g\otimes_\pi R\) 映入
\(U(\mathfrak g)\) 的中心。对于 \(x\in\mathfrak g\)，令
\[
  x^{(p)}
  =
  x^p-\gamma(x\otimes_\pi1).
\]

% 权威底本： Exp7A 局部页 35; 合并本页 488; 印刷页 477.
则对于每个 \(\alpha\) 都有 \(x_\alpha^{(p)}=y_\alpha\)。若
\[
  x=\sum_\alpha\lambda_\alpha x_\alpha,
\]
\SGARefLink{sga3:VIIA:5.1}{公式 5.1}~\SGARefLink{sga3:vii:paragraph:5-1:equation:ii}{(ii)}，对满足 \(\lambda_\alpha\ne0\) 的指标 \(\alpha\) 的个数作归纳，得到
\[
  x^p-\sum_\alpha\lambda_\alpha^p x_\alpha^p
  \in\mathfrak g.
\]
将此元素记为 \(z\)。则
\[
  x^{(p)}
  =
  \sum_\alpha\lambda_\alpha^p y_\alpha+z,
\]
故 \(x^{(p)}\in\mathfrak g\)。

显然
\[
  (\lambda x)^{(p)}
  =
  \lambda^p x^{(p)}.
\]
由于 \(\gamma(x\otimes_\pi1)\) 位于中心，有
\(\operatorname{ad}x^{(p)}=\operatorname{ad}x^p\)。因此 Jacobson 的第一个公式 \SGARefLink{sga3:VIIA:5.1}{5.1}~\SGARefLink{sga3:vii:paragraph:5-1:equation:i}{(i)} 给出
\[
  \operatorname{ad}x^{(p)}
  =
  (\operatorname{ad}x)^p.
\]
最后，Jacobson 的第二个公式 \SGARefLink{sga3:VIIA:5.1}{5.1}~\SGARefLink{sga3:vii:paragraph:5-1:equation:ii}{(ii)} 表明映射
\(x\mapsto x^{(p)}\) 满足 \SGARefLink{sga3:vii:paragraph:5-1:equation:ii}{(ii)} （见 \SGARefLink{sga3:VIIB:5.2:theorem}{5.2}）的条件。因此它使 \(\mathfrak g\) 成为受限 \(p\)-Lie 代数。命题得证。

\medskip
\noindent\textbf{5.3.3. 命题。}
\phantomsection
\label{sga3:VIIA:5.3.3}
\textit{设 \(\mathfrak g\) 是 \(R\) 上的受限 \(p\)-Lie 代数，其底层模为自由模，基为 \((x_\alpha)\)。则
\[
  j:\mathfrak g\longrightarrow U_p(\mathfrak g)
\]
是单射。若 \(z_\alpha=j(x_\alpha)\)，则 \(U_p(\mathfrak g)\) 以如下单项式为基：
\[
  \prod_\alpha z_\alpha^{n_\alpha},
  \qquad
  0\le n_\alpha<p.
\]
这里除有限多个 \(\alpha\) 外，\(n_\alpha\) 均为零；假设基带有全序，乘积按递增顺序取。}

通过典范映射 \(i\)，将 \(\mathfrak g\) 等同于包络代数 \(U(\mathfrak g)\) 的一个子模。对于每个除有限多项外均为零的非负整数族
\[
  n=(n_\alpha)
\]
，令
\[
  |n|=\sum_\alpha n_\alpha,
  \qquad
  x^n=\prod_\alpha x_\alpha^{n_\alpha}.
\]
写成
\[
  n_\alpha=m_\alpha+p\ell_\alpha,
  \qquad
  0\le m_\alpha<p,
\]
并令
\[
  T_n
  =
  \prod_\alpha
  x_\alpha^{m_\alpha}
  \beta(x_\alpha)^{\ell_\alpha},
\]
其中
\[
  \beta(x)=x^p-x^{(p)}
\]
是 \SGARefLink{sga3:VIIA:5.3.1}{5.3.1} 中定义的映射 \(\mathfrak g\to U(\mathfrak g)\)。

对于 \(r\in\mathbb N\)，令 \(U^r\) 为 \(U(\mathfrak g)\) 中由满足 \(|n|\le r\) 的 \(x^n\) 生成的 \(R\)-子模。由于分次环
\[
  \bigoplus_r U^r/U^{r-1}
\]
是交换的（参见 \textup{[BLie]}, I, \S2.6），对于每个 \(n\) 都有
\[
  T_n-\prod_\alpha x_\alpha^{n_\alpha}
  \in U^{|n|-1}.
\]
对于每个 \(s\in\mathbb N\)，由 Poincar\'e--Birkhoff--Witt 定理（\emph{同上}，\S2.7），满足 \(|n|=s\) 的 \(x^n\) 构成 \(U^s/U^{s-1}\) 的一组基；因此满足 \(|n|=s\) 的 \(T_n\) 也如此。

因而，当 \(s=|n|\) 变化时，\(T_n\) 构成 \(U(\mathfrak g)\) 的一组基。由 \SGARefLink{sga3:VIIA:5.3.1}{5.3.1}，典范映射
\[
  U(\mathfrak g)\longrightarrow U_p(\mathfrak g)
\]
的核 \(J\) 是 \(U(\mathfrak g)\) 中由中心元素 \(\beta(x_\alpha)\) 生成的左理想。因此，满足 \(\ell=(\ell_\alpha)\ne0\) 的 \(T_n\) 构成 \(J\) 的一组基，而对每个 \(\alpha\) 均满足 \(n_\alpha<p\) 的 \(T_n\) 构成下式的一组基：
\[
  U_p(\mathfrak g)=U(\mathfrak g)/J.
\]

\medskip
\noindent\textbf{5.3.3 bis.}
\phantomsection
\label{sga3:VIIA:5.3.3bis}

% 权威底本： Exp7A 局部页 36; 合并本页 489; 印刷页 478.
设 \(\mathfrak g\) 是 \(R\) 上的受限 \(p\)-Lie 代数，\(f:R\to R'\) 是基环的扩张。我们断言，\(R'\)-模
\[
  R'\otimes_R\mathfrak g
\]
容许唯一的受限 \(p\)-Lie 代数结构，使得
\[
  \tag{*}\label{sga3:vii:paragraph:5-3-3-bis:equation:star}
  [\lambda\otimes x,\mu\otimes y]
  =
  \lambda\mu\otimes[x,y],
  \qquad
  (\lambda\otimes x)^{(p)}
  =
  \lambda^p\otimes x^{(p)}.
\]
特别地，由此可知函子
\[
  \mathfrak g\longmapsto R'\otimes_R\mathfrak g
\]
是从 \(R'\) 到 \(R\) 的标量限制的左伴随。

由 \((*)\) 定义的结构的唯一性是显然的；下面证明其存在性。若 \(\mathfrak g\) 是以 \((x_\alpha)\) 为基的自由模，则 \SGARefLink{sga3:VIIA:5.3.2}{命题 5.3.2} 给出 Lie 代数 \(R'\otimes_R\mathfrak g\) 上唯一的受限 \(p\)-Lie 代数结构，使得
\[
  (1\otimes x_\alpha)^{(p)}
  =
  1\otimes x_\alpha^{(p)}.
\]
这就是所求的结构。

现在设 \(\mathfrak g\) 是任意的。存在一个作为 \(R\)-模为自由模的受限 \(p\)-Lie 代数 \(L_0\)，以及满同态
\[
  q_0:L_0\longrightarrow\mathfrak g.
\]
例如，可以取
\[
  L_0=R\otimes_{\mathbb F_p}\mathfrak g
\]
并令 \(q_0\) 为同态
\[
  \lambda\otimes x\longmapsto\lambda x;
\]
这里 \(\mathfrak g\) 在素域 \(\mathbb F_p\) 上是自由的。\(q_0\) 的核是 \(L_0\) 的一个 \(p\)-理想，即 Lie 代数 \(L_0\) 中在 \(x\mapsto x^{(p)}\) 下稳定的理想。因此还存在一个作为 \(R\)-模为自由模的受限 \(p\)-Lie 代数 \(L_1\)，以及像为 \(\operatorname{Ker}q_0\) 的同态 \(q_1:L_1\to L_0\)。这给出正合列
\begin{center}
  \phantomsection
  \label{sga3:VIIA:5.3.3bis-exact-sequence}
  % SGA3-VII-DIAG-050
  \(\displaystyle
L_1 \xrightarrow{q_1} L_0 \xrightarrow{q_0} \mathfrak g \longrightarrow 0 .
\)

\end{center}
因而给出 \(R'\)-Lie 代数的正合列
\begin{center}
  \phantomsection
  \label{sga3:VIIA:5.3.3bis-basechange-exact-sequence}
  % SGA3-VII-DIAG-051
  \(\displaystyle
R'\otimes_R L_1
  \xrightarrow{R'\otimes_R q_1}
R'\otimes_R L_0
  \xrightarrow{R'\otimes_R q_0}
R'\otimes_R\mathfrak g
  \longrightarrow 0 .
\)

\end{center}

由于 \(R'\otimes_Rq_1\) 显然是受限 \(p\)-Lie 代数同态，\(R'\otimes_Rq_0\) 的核是一个 \(p\)-理想。因此，\(R'\otimes_RL_0\) 上的符号 \(p\) 次幂运算通过取商，诱导出 \(R'\otimes_R\mathfrak g\) 到自身的映射（使用 \SGARefLink{sga3:VIIB:5.2:theorem}{5.2} 的公式 \SGARefLink{sga3:vii:paragraph:5-1:equation:ii}{(ii)}）。此映射赋予 \(R'\otimes_R\mathfrak g\) 所需的受限 \(p\)-Lie 代数结构。

\medskip
\noindent\textbf{5.3.4.}
\phantomsection
\label{sga3:VIIA:5.3.4}

典范映射
\[
  j_{\mathfrak g}:\mathfrak g\longrightarrow U_p(\mathfrak g)
\]
对于基环的每个扩张 \(R\to R'\)，诱导出同态
\[
  R'\otimes_Rj_{\mathfrak g}:
  R'\otimes_R\mathfrak g
  \longrightarrow
  R'\otimes_RU_p(\mathfrak g).
\]
因此它诱导出同态
\[
  h:
  U_p(R'\otimes_R\mathfrak g)
  \longrightarrow
  R'\otimes_RU_p(\mathfrak g)
\]
使得
\[
  h\circ j_{R'\otimes_R\mathfrak g}
  =
  R'\otimes_Rj_{\mathfrak g}.
\]
由 \(R'\otimes_R\mathfrak g\) 和受限包络代数的泛性质，立即可知 \(h\) 是同构。因此可作如下等同：
\[
  U_p(R'\otimes_R\mathfrak g)
  \simeq
  R'\otimes_RU_p(\mathfrak g).
\]

% 权威底本： Exp7A 局部页 37; 合并本页 490; 印刷页 479.
特别地，设 \(r\in R\)，令 \(R'=R_r\) 为局部化环。则
\[
  \mathfrak g_r=R_r\otimes_R\mathfrak g
\]
具有典范的 \(R_r\) 上受限 \(p\)-Lie 代数结构。因而，\(\operatorname{Spec}R\) 上的层 \(\widetilde{\mathfrak g}\) 是 \(\operatorname{Spec}R\) 上的拟凝聚受限 \(p\)-Lie 代数。此外，
\[
  U_p^{R_r}(\mathfrak g_r)
  \simeq
  U_p^R(\mathfrak g)_r,
\]
因此，与预层
\[
  V\longmapsto
  U_p\bigl(\Gamma(V,\widetilde{\mathfrak g})\bigr)
\]
相伴的层是拟凝聚的。

\noindent\textbf{定义。}
更一般地，设 \(S\) 是特征为 \(p\) 的概形，\(\mathcal G\) 是 \(\mathcal O_S\) 上的拟凝聚受限 \(p\)-Lie 代数。与预层
\[
  V\longmapsto U_p\bigl(\Gamma(V,\mathcal G)\bigr)
\]
相伴的层是拟凝聚的。将其记为 \(\mathcal U_p(\mathcal G)\)，并称为 \(\mathcal G\) 的\emph{受限包络代数}。若 \(V\) 是仿射的，则
\[
  U_p\bigl(\Gamma(V,\mathcal G)\bigr)
  \simeq
  \Gamma\bigl(V,\mathcal U_p(\mathcal G)\bigr).
\]

% 下一游标： Exp7A 局部页 37, 合并阅读本页 490,
% 再版印刷页 479, 位于 5.4。
